\section*{Цель работы}
Изучение зависимости между деформациями и напряжениями при осевом растяжении металлического образца, а также определение механических характеристик прочности и пластичности материала.

\section*{Краткие теоретические сведения}
Под действием внешних сил твердое тело деформируется. Различают упругие деформации (исчезающие после снятия нагрузки) и пластические (остаточные).
Зависимость между нагрузкой $P$ и абсолютным удлинением $\Delta l$ описывается диаграммой растяжения.

Основные характерные точки диаграммы и соответствующие им напряжения:
\begin{enumerate}
    \item \textit{Предел пропорциональности} ($\sigma_{\text{пц}}$) — напряжение, до которого справедлив закон Гука:
    $$ \sigma = E \cdot \varepsilon $$
    где $E$ — модуль упругости (Юнга), $\varepsilon$ — относительная деформация.
    \item \textit{Предел текучести} ($\sigma_{\text{т}}$) — напряжение, при котором деформация растет без заметного увеличения нагрузки (площадка текучести). При отсутствии выраженной площадки используют условный предел текучести $\sigma_{0.2}$ (напряжение, вызывающее остаточную деформацию $0.2\,\%$).
    \item \textit{Временное сопротивление} ($\sigma_{\text{в}}$) — напряжение, соответствующее максимальной нагрузке $P_{\max}$, которую выдерживает образец до разрушения.
    \item \textit{Истинное сопротивление разрыву} ($\sigma_{\text{к}}$) — напряжение в момент разрыва, определяемое как отношение нагрузки $P_{\text{к}}$ к площади шейки $F_{\text{к}}$.
\end{enumerate}

Для оценки пластичности используют:
\begin{itemize}
    \item \textit{Относительное удлинение после разрыва} ($\delta$):
    $$ \delta = \frac{l_{\text{к}} - l_0}{l_0} \cdot 100\,\% $$
    \item \textit{Относительное сужение после разрыва} ($\psi$):
    $$ \psi = \frac{F_0 - F_{\text{к}}}{F_0} \cdot 100\,\% $$
\end{itemize}

\section*{Оборудование и материалы}
\begin{itemize}
    \item Испытательная установка (разрывная машина) с системой управления и ПК.
    \item Штангенциркуль.
    \item Образец для испытаний (стандартный цилиндрический).
\end{itemize}

\begin{figure}[H]
    \centering
    \includegraphics[width=0.6\linewidth]{figs/machine_scheme.pdf}
    \caption{Блок-схема испытательной установки}
    \label{fig:setup}
\end{figure}

\section*{Инструкция по снятию данных}
Процесс снятия экспериментальных данных является критическим этапом. Ошибки в измерениях геометрических размеров делают дальнейшие расчеты бессмысленными.

\begin{enumerate}
    \item \textit{Подготовка образца и начальные измерения.}
    До установки образца в захваты машины:
    \begin{itemize}
        \item Измерьте начальную длину рабочей части $l_0$ штангенциркулем.
        \item Измерьте начальный диаметр $d_0$ в трех сечениях (по краям и в центре рабочей части). В каждом сечении проведите измерения в двух взаимно перпендикулярных направлениях.
        \item Занесите наименьшее среднее значение диаметра в протокол.
    \end{itemize}

    \item \textit{Установка образца.}
    \begin{itemize}
        \item Закрепите образец в захватах разрывной машины.
        \item Убедитесь в отсутствии перекосов. Образец должен быть центрирован.
        \item Обнулите показания датчиков силы и перемещения на пульте управления или в ПО.
    \end{itemize}

    \item \textit{Проведение испытания.}
    \begin{itemize}
        \item Запустите процесс нагружения с заданной скоростью (согласно указаниям преподавателя или лаборанта).
        \item Наблюдайте за построением диаграммы растяжения на мониторе ПК.
        \item Дождитесь полного разрушения образца.
        \item После разрыва остановите машину (если автоматическая остановка не сработала).
    \end{itemize}

    \item \textit{Фиксация нагрузок.}
    Используя программное обеспечение или распечатку диаграммы, определите и занесите в протокол значения нагрузок:
    \begin{itemize}
        \item $P_{\text{пц}}$ (нагрузка предела пропорциональности);
        \item $P_{\text{т}}$ (нагрузка предела текучести);
        \item $P_{\max}$ (максимальная нагрузка);
        \item $P_{\text{к}}$ (нагрузка в момент разрыва).
    \end{itemize}

    \item \textit{Финальные измерения.}
    \begin{itemize}
        \item Извлеките разрушенные части образца из захватов.
        \item Плотно состыкуйте части образца в месте разрыва.
        \item Измерьте конечную длину рабочей части $l_{\text{к}}$.
        \item Измерьте минимальный диаметр шейки $d_{\text{к}}$ в месте разрыва (в двух взаимно перпендикулярных направлениях).
        \item Занесите данные в протокол.
    \end{itemize}
\end{enumerate}

\section*{Инструкция по обработке результатов}
Обработка выполняется в файле \texttt{processing.tex} и включает следующие этапы:
\begin{enumerate}
    \item Вычисление площадей поперечного сечения до ($F_0$) и после ($F_{\text{к}}$) испытания.
    \item Расчет механических характеристик прочности ($\sigma_{\text{пц}}, \sigma_{\text{т}}, \sigma_{\text{в}}, \sigma_{\text{к}}$).
    \item Расчет характеристик пластичности ($\delta, \psi$).
    \item Построение диаграммы условных напряжений $\sigma = f(\varepsilon)$.
    \item Построение диаграммы истинных напряжений $s = f(e)$ (дополнительное задание).
    \item Формулировка вывода о типе разрушения (вязкое/хрупкое) и соответствии материала классу сталей.
\end{enumerate}
